%% Documento de arquitectura del sistema
%% Compilar: tectonic -X compile documento-arquitectura.tex
\documentclass[11pt,a4paper]{article}
\usepackage{../latex/legasoft}

\lgtitulo{Documento de arquitectura del sistema}
\lgsubtitulo{Descripción técnica basada en el modelo C4}
\lgtituloCorto{Arquitectura}
\lgcodigo{LGS-ARQ-001}
\lgversion{0.1}
\lgestado{Plantilla}
\lgfecha{\campo{fecha}}
\lgautor{\lgArquitecto\ (\lgArquitectoRol)}

\begin{document}
\lgportada

\lgcontrolversiones{
0.1 & \campo{dd-mm-aaaa} & \lgArquitecto & Plantilla inicial basada en el modelo C4. \\
}

\begin{porcompletar}[Cómo usar esta plantilla]
La representación de la arquitectura mediante el modelo C4 es responsabilidad del
\lgArquitectoRol\ (\texttt{LGS-INST-001}, sección 10.2). Los diagramas fuente
viven en \ruta{modelos/c4/} y se incorporan aquí como figuras una vez
compilados. Mientras no existan decisiones tomadas, esta plantilla documenta la
estructura esperada, no una arquitectura definida.
\end{porcompletar}

\tableofcontents
\clearpage

% =============================================================================
\section{Introducción}

\subsection{Propósito}

Este documento describe la arquitectura del sistema \campo{nombre del sistema}:
su estructura, las decisiones técnicas que la sostienen y las consecuencias de
esas decisiones para la evolución del producto. Está dirigido al equipo de
desarrollo y a quien deba mantener la solución después de la entrega.

\subsection{Alcance}

\begin{porcompletar}
\campo{Partes del sistema cubiertas por esta descripción y aquellas que quedan
fuera.}
\end{porcompletar}

\subsection{Documentos relacionados}

\begin{itemize}
  \item \texttt{LGS-REQ-001} — Especificación de requisitos de software.
  \item \texttt{LGS-CAL-001} — Estrategia de calidad y plan de pruebas.
  \item \ruta{modelos/c4/} — Diagramas de contexto, contenedores y componentes.
  \item \ruta{docs/arquitectura/adr/} — Registros de decisiones de arquitectura.
\end{itemize}

% =============================================================================
\section{Restricciones y objetivos de calidad}

La arquitectura debe ser proporcional al problema. Legasoft evita introducir
complejidad sin un beneficio verificable, y ese criterio gobierna la selección de
plataformas, patrones, servicios en la nube y mecanismos de automatización.

\subsection{Restricciones}

\begin{porcompletar}
\campo{Tecnologías impuestas, infraestructura disponible, presupuesto, plazos,
normativa o capacidad operativa del cliente.}
\end{porcompletar}

\subsection{Atributos de calidad prioritarios}

Los atributos se toman de \texttt{LGS-REQ-001} y se expresan aquí como fuerzas que
condicionan el diseño.

\begin{center}
\begin{tabularx}{\linewidth}{@{}L{3.4cm} C{1.8cm} Y@{}}
\toprule
\sffamily\bfseries\color{lgPrimario}Atributo &
\sffamily\bfseries\color{lgPrimario}Prioridad &
\sffamily\bfseries\color{lgPrimario}Implicación para la arquitectura \\
\midrule
Mantenibilidad & \campo{A/M/B} & \campo{modularidad, separación de responsabilidades.} \\
Seguridad      & \campo{A/M/B} & \campo{autenticación, autorización, cifrado, auditoría.} \\
Escalabilidad  & \campo{A/M/B} & \campo{estrategia de crecimiento previsto.} \\
Fiabilidad     & \campo{A/M/B} & \campo{manejo de fallos y recuperación.} \\
Desempeño      & \campo{A/M/B} & \campo{tiempos de respuesta y volumen esperado.} \\
\bottomrule
\end{tabularx}
\captionof{table}{Atributos de calidad que condicionan el diseño.}
\end{center}

% =============================================================================
\section{Vista de contexto (C4 nivel 1)}

Muestra el sistema como una caja negra y su relación con usuarios y sistemas
externos.

\begin{porcompletar}[Diagrama pendiente]
Compilar \ruta{modelos/c4/c4-contexto.tex} e incorporar el resultado aquí:

\vspace{4pt}
\texttt{\textbackslash includegraphics[width=\textbackslash linewidth]\{../../modelos/c4/c4-contexto.pdf\}}
\end{porcompletar}

\begin{porcompletar}
\campo{Descripción de los actores, los sistemas externos y el propósito de cada
relación.}
\end{porcompletar}

% =============================================================================
\section{Vista de contenedores (C4 nivel 2)}

Describe las unidades desplegables del sistema, es decir las aplicaciones, los
servicios y los almacenes de datos, y cómo se comunican entre sí.

\begin{porcompletar}[Diagrama pendiente]
Compilar \ruta{modelos/c4/c4-contenedores.tex} e incorporar el resultado aquí.
\end{porcompletar}

\begin{center}
\begin{tabularx}{\linewidth}{@{}L{3.4cm} L{3.2cm} Y@{}}
\toprule
\sffamily\bfseries\color{lgPrimario}Contenedor &
\sffamily\bfseries\color{lgPrimario}Tecnología &
\sffamily\bfseries\color{lgPrimario}Responsabilidad \\
\midrule
\campo{Aplicación web} & \campo{por definir} & \campo{responsabilidad} \\
\campo{API}            & \campo{por definir} & \campo{responsabilidad} \\
\campo{Base de datos}  & \campo{por definir} & \campo{responsabilidad} \\
\bottomrule
\end{tabularx}
\captionof{table}{Contenedores del sistema y su responsabilidad.}
\end{center}

\begin{nota}[Decisiones tecnológicas]
Cada tecnología de frontend, backend y pruebas se justifica por las necesidades,
restricciones y recursos del proyecto, y se registra como ADR en
\ruta{docs/arquitectura/adr/}.
\end{nota}

% =============================================================================
\section{Vista de componentes (C4 nivel 3)}

Detalla la estructura interna de los contenedores relevantes.

\begin{porcompletar}[Diagrama pendiente]
Compilar \ruta{modelos/c4/c4-componentes.tex} e incorporar el resultado aquí.
\end{porcompletar}

% =============================================================================
\section{Modelo de datos}

\begin{porcompletar}
\campo{Entidades principales, relaciones, reglas de integridad y decisiones de
modelado. Cuando el contexto lo justifique podrá emplearse SQL Server como gestor
relacional.}
\end{porcompletar}

% =============================================================================
\section{Contratos de integración}

Los contratos entre frontend y backend se documentan para reducir ambigüedades y
permitir que el \lgQARol\ diseñe pruebas de integración sobre comportamientos
acordados.

\begin{center}
\begin{tabularx}{\linewidth}{@{}L{1.9cm} L{3.0cm} Y C{2.7cm}@{}}
\toprule
\sffamily\bfseries\color{lgPrimario}Método &
\sffamily\bfseries\color{lgPrimario}Ruta &
\sffamily\bfseries\color{lgPrimario}Propósito &
\sffamily\bfseries\color{lgPrimario}Autorización \\
\midrule
\campo{GET}  & \campo{/recurso}      & \campo{descripción} & \campo{rol} \\
\campo{POST} & \campo{/recurso}      & \campo{descripción} & \campo{rol} \\
\bottomrule
\end{tabularx}
\captionof{table}{Resumen de los contratos expuestos por la API.}
\end{center}

% =============================================================================
\section{Seguridad}

\begin{porcompletar}
\campo{Mecanismos de autenticación y autorización, tratamiento de datos
sensibles, registro de auditoría y controles sobre acceso, confidencialidad,
integridad y disponibilidad, tomando ISO/IEC 27001 como referencia.}
\end{porcompletar}

% =============================================================================
\section{Despliegue y operación}

La adopción de prácticas de Cloud y DevOps es gradual y debe considerar
seguridad, costos, capacidad operativa y continuidad. Su propósito es reducir
errores manuales, mejorar la repetibilidad de las entregas y conservar evidencia
sobre los cambios realizados.

\begin{porcompletar}
\campo{Entornos previstos, estrategia de despliegue, contenerización, pipeline de
integración y entrega continua, y mecanismos de monitoreo. Los workflows se
alojarán en .github/workflows/.}
\end{porcompletar}

% =============================================================================
\section{Decisiones de arquitectura}

Cada decisión relevante se registra como un ADR independiente, con su contexto,
las alternativas evaluadas y sus consecuencias. La plantilla se encuentra en
\ruta{docs/arquitectura/adr-plantilla.tex}.

\begin{center}
\begin{tabularx}{\linewidth}{@{}C{2.0cm} Y C{2.4cm} C{2.2cm}@{}}
\toprule
\sffamily\bfseries\color{lgPrimario}ADR &
\sffamily\bfseries\color{lgPrimario}Decisión &
\sffamily\bfseries\color{lgPrimario}Estado &
\sffamily\bfseries\color{lgPrimario}Fecha \\
\midrule
ADR-001 & \campo{decisión} & \campo{propuesta} & \campo{fecha} \\
ADR-002 & \campo{decisión} & \campo{propuesta} & \campo{fecha} \\
\bottomrule
\end{tabularx}
\captionof{table}{Índice de decisiones de arquitectura registradas.}
\end{center}

% =============================================================================
\section{Riesgos técnicos y deuda conocida}

\begin{porcompletar}
\campo{Riesgos técnicos identificados, deuda asumida conscientemente y condiciones
en las que debería revisarse.}
\end{porcompletar}

\end{document}
